\subsection{Sample flow and recruitment}
\label{sec:results-flow}
Table~\ref{tab:sample-flow} summarizes the flow from survey completion to the balanced analysis cohort of 352 athletes and 704 athlete-period records. Phase 2 targeted recall concentrated additional enrollment in Under 20 and Under 23 strata, particularly hurdles, horizontal jumps, and combined events, which the voluntary Phase 1 sample under-represented relative to the FIDAL population estimate. After poststratification (Section~\ref{sec:weighting}; Appendix Table~\ref{tab:weight-diagnostics}), the Kish effective sample size was 340.5 of 352 analyzed athletes (97\%), indicating that reweighting did not substantially inflate posterior uncertainty.

\subsection{Baseline characteristics}
\label{sec:results-table1}
Table~\ref{tab:table1} compares baseline characteristics by competition sex, the principal descriptive contrast for this cohort, alongside the corresponding population-weighted standardized mean differences (SMD). As expected for a sex comparison, height and weight showed large SMDs ($-1.81$ and $-1.83$); demographic, injury-history, carbon-exposure, and discipline-distribution SMDs were mostly below the conventional 0.20 threshold for a small imbalance \citep{table1guidelines}, with horizontal jumps and combined events showing modest over-representation of women (SMD 0.22--0.23). Reweighting to national population strata (weighted SMD column) did not materially change this pattern, supporting the sample as broadly representative of the target competitive population once poststratified.

\subsection{Descriptive carbon-plate insole use}
\label{sec:results-descriptive}
Across the balanced cohort, carbon-plated spikes accounted for a mean 35.7\% of reported preparation-period training-type frequency and 47.7\% of reported competition-period training-type frequency (Table~\ref{tab:descriptive-carbon}). This pattern--roughly two-fifths of preparation training and roughly half of competition training reported in carbon-plated spikes--was stable across sex and age-group strata (Table~\ref{tab:descriptive-carbon}, Figure~\ref{fig:descriptive-carbon-share-by-stratum}), with mean shares differing by at most 4 percentage points between men and women or across the three age groups in either period, and with the full distributions--not only their means--overlapping closely between strata. Adoption (any carbon use) was high in both periods: 81.5\% of athletes reported any preparation carbon use and 87.5\% reported any competition carbon use (Table~\ref{tab:table1}), consistent with carbon-plated spikes functioning as the prevailing rather than niche equipment choice in this population.

\subsection{Specialty-group-stratified characteristics}
\label{sec:results-discipline}
To summarize the seven disciplines with limited individual cell sizes (19--107 athletes; Appendix Table~\ref{tab:discipline-support}), they were collapsed into four broader, physiologically motivated specialty groups--sprint/hurdles, endurance (middle and long distance), jumps (horizontal jumps or pole vault), and combined events. Tables~\ref{tab:specialty-group} and~\ref{tab:specialty-group-carbon} report baseline characteristics, injury prevalence, and carbon-plate use by group. Because disciplines and, consequently, groups are multi-select, 26 of 352 athletes belong to more than one group. Two patterns stand out as original descriptive findings of this cohort.

First, competition-period carbon share tracked the explosive-versus-endurance character of the group rather than being uniform across athletics. The sprint/hurdles group (the largest, $n=169$) had the highest competition carbon share (0.56) and the largest preparation-to-competition rise (+0.17), while the endurance group ($n=109$) had the lowest carbon share in both periods (0.31 and 0.34). This is consistent with the biomechanical rationale in Section~\ref{sec:introduction}: carbon-plate stiffness benefits are most established for short-duration, high-power efforts, and this cohort's self-reported equipment allocation mirrors that distinction even though the survey did not ask athletes about performance rationale directly.

Second, injury prevalence did not track carbon share in the same way. The jumps group ($n=62$) stood out with the highest prior-injury history (83.9\%) and the highest competition-period injury prevalence (51.6\%) of the four groups, despite an intermediate carbon share (0.47)--higher than endurance but lower than sprint/hurdles--indicating that the elevated injury burden in jumping disciplines is unlikely to be explained by carbon-plate dose alone and more plausibly reflects the landing and eccentric loading demands specific to jumping events; the primary Bayesian model (Section~\ref{sec:results-primary}) already adjusts for discipline using partially pooled indicators specifically to avoid attributing this discipline-driven variation to carbon exposure. The combined-events group ($n=41$) showed a distinct hybrid profile: a carbon share close to sprint/hurdles (0.52) paired with an endurance-like training volume (preparation and competition medians of 9) and an elevated any-injury prevalence (70.7\%), consistent with the mixed physiological demands of multi-event competition and with the higher, more heterogeneous musculoskeletal load noted for plyometric and multi-skill disciplines in prior athletics injury literature \citep{edouard2020injury}.

\subsection{Injury location and type}
\label{sec:results-pattern}
The 204 injured athletes reported 259 injuries in total (149 athletes with one injury, 55 with two), a median of 1 injury per affected athlete. Table~\ref{tab:injury-location} reports the anatomical location of these injuries and Figure~\ref{fig:injury-location-type} displays both the location and type distributions together. The thigh was, by a substantial margin, the single most commonly reported location (97 of 259 injuries, 37.5\%; affecting 80 of 204 injured athletes, 39.2\%), followed by the lower leg (25.1\% of injuries) and the foot (19.3\%); together, the hip, groin, thigh, lower leg, ankle, and foot accounted for 95.8\% of all reported injuries, with only 4.2\% localized to the head, upper limb, or trunk.

For the 248 lower-limb injuries for which an injury type was recorded (Table~\ref{tab:injury-type}, Figure~\ref{fig:injury-location-type}), tendinopathy (28.6\%) and bone stress injury or stress fracture (26.2\%) were the two most frequently reported types, together accounting for more than half of classified lower-limb injuries and each individually more common than muscle strain (22.2\%). This is a distinctive feature of this cohort relative to general athletics injury surveillance, where muscle injury is typically the most frequently reported tissue type by a wide margin and tendon and bone injuries are reported considerably less often \citep{edouard2020injury}.

Cross-tabulating type against location (Table~\ref{tab:injury-type-by-location}) shows that these two leading types were not distributed the same way anatomically. Tendinopathy was disproportionately distal relative to the overall injury-location pattern: foot (35.2\% of the 71 tendinopathy injuries) and lower leg (33.8\%) were more prominent locations for tendinopathy than for the injury sample as a whole (19.3\% and 25.1\% respectively; Table~\ref{tab:injury-location}), and ankle involvement (22.5\%) was likewise higher than the sample-wide 13.9\%--a pattern compatible with Achilles- and plantar-tendon structures at the distal lower limb, consistent with the biomechanical literature on carbon-plate loading of the ankle and forefoot \citep{gerlach2025aft}. Bone stress injury or stress fracture, in contrast, was thigh-dominant (58.5\% of the 65 bone stress injuries), a pattern closely tracking the overall injury-location distribution and similar to muscle strain (74.5\% thigh) rather than being disproportionately distal; foot involvement among bone stress injuries (15.4\%) was, if anything, lower than the sample-wide rate. The elevated overall frequency of bone stress injury in this cohort is therefore a genuine descriptive finding, but this cohort's data do not show the foot/navicular-specific concentration that would most directly support the mechanistic link to carbon-plate loading proposed for that injury type \citep{tenforde2023bonestress}; the survey's type item does not further subclassify bone site (for example, navicular versus tibial versus metatarsal), so this distinction cannot be resolved further with the present data. Taken together, the location-by-type pattern strengthens the tendinopathy--carbon-plate hypothesis specifically, while leaving the bone-stress-injury hypothesis only partially supported. This is reported as a hypothesis-generating observation: the present cross-sectional design and sample size do not support a formal type-stratified test of the carbon-share association, and confirming a tendinopathy-specific association with carbon-plate dose, or resolving the anatomical site of the bone stress injuries observed here, would require a dedicated, adequately powered follow-up.

\subsection{Athlete-reported injury burden}
\label{sec:results-ostrc}
Of the 204 injured athletes, 193 (94.6\%) completed all four adapted OSTRC-H2 impact items. The median total severity score was 75 (IQR 58--92) overall (Table~\ref{tab:ostrc-period}). Scores were similar for athletes injured only during preparation (median 71.5, IQR 50--90) and only during competition (67, IQR 58--84; SMD $-0.07$), while athletes reporting injury in both periods had a higher median score of 84 (IQR 67--92; SMD 0.41 versus competition only). The largest component difference was for symptoms in the both-period group (SMD 0.51). These results describe impact among athletes already classified as injured; they are not a second injury-prevalence analysis.

\begin{table}[htbp]
\centering
\small
\caption{Adapted OSTRC-H2 impact scores among 193 injured athletes with complete responses, stratified by mutually exclusive injury period. Values are median [Q1, Q3]. Each component ranges from 0 to 25 and total severity from 0 to 100; higher scores indicate greater impact. SMDs compare group means with the competition-only group and are descriptive.}
\label{tab:ostrc-period}
\begin{tabular}{L{2.8cm}rrrrrr}
\toprule
& Overall & Preparation & Competition & Both & \multicolumn{2}{c}{SMD vs competition} \\
\cmidrule(lr){6-7}
Domain & ($n=193$) & only ($n=50$) & only ($n=84$) & periods ($n=59$) & Preparation & Both \\
\midrule
Participation & 17 [17, 25] & 17 [17, 25] & 17 [8, 25] & 17 [17, 25] & 0.14 & 0.19 \\
Training volume & 17 [17, 25] & 17 [8, 25] & 17 [8, 25] & 25 [17, 25] & 0.01 & 0.29 \\
Performance & 17 [17, 25] & 17 [8, 25] & 17 [17, 25] & 25 [17, 25] & $-0.35$ & 0.22 \\
Symptoms & 17 [8, 25] & 17 [8, 17] & 17 [8, 17] & 17 [17, 25] & $-0.01$ & 0.51 \\
Total severity & 75 [58, 92] & 71.5 [50, 90] & 67 [58, 84] & 84 [67, 92] & $-0.07$ & 0.41 \\
\bottomrule
\end{tabular}
\end{table}


\begin{figure}[htbp]
\centering
\includegraphics[width=0.95\textwidth]{fig_injury_location_type.png}
\caption{Anatomical location (\% of 259 injuries) and type (\% of 248 lower-limb injuries with a recorded type) of injuries reported by the balanced analysis cohort. Both items are multi-select, so percentages within each panel sum to more than 100\%.}
\label{fig:injury-location-type}
\end{figure}

\subsection{Primary adjusted associations}
\label{sec:results-primary}
Figure~\ref{fig:carbon-share-distribution} shows the observed distribution of carbon share by period; Figure~\ref{fig:spline-dose-response} displays the fitted among-user spline, and Figure~\ref{fig:posterior-associations} and Table~\ref{tab:primary-associations} summarize the primary adjusted estimands. Adjusted preparation-period injury prevalence was 30.1\% (95\% CrI 21.0--40.2\%) for non-users and 35.5\% (95\% CrI 29.6--41.7\%) for median-share users; the adoption difference of 5.4 percentage points (pp) (95\% CrI $-6.3$ to $16.2$) had a posterior probability of 82.0\% of being positive, well short of the threshold used here to describe an association as clear. The preparation-period dose contrast (75th vs.\ 25th percentile share among users) was $-0.8$ pp (95\% CrI $-10.6$ to $9.1$; posterior probability of a positive association 43.8\%), that is, centered near no association.

Adjusted competition-period injury prevalence was 38.1\% (95\% CrI 26.6--50.4\%) for non-users and 43.0\% (95\% CrI 36.5--49.4\%) for median-share users, an adoption difference of 4.8 pp (95\% CrI $-8.7$ to $17.7$; posterior probability 77.1\%). In contrast, among competition-period users, moving from the 25th percentile (36.0\% share) to the 75th percentile (71.4\% share) of carbon share was associated with adjusted prevalence rising from 36.0\% (95\% CrI 28.5--44.0\%) to 50.3\% (95\% CrI 42.4--58.2\%): a difference of 14.3 pp (95\% CrI 2.7 to 25.8) and a ratio of 1.41 (95\% CrI 1.06 to 1.86), with a posterior probability of 99.0\% that the association was positive and 94.0\% that it exceeded the pre-specified 5-pp threshold for a practically relevant association (Table~\ref{tab:primary-associations}). This competition-period dose association was the only one of the four primary contrasts for which the credible interval excluded zero.

\subsection{Model validation}
\label{sec:results-validation}
Both the primary (exposure) and matched baseline (no-exposure) models showed acceptable sampling behavior: maximum $\hat{R} \leq 1.003$, zero divergent transitions, zero maximum-tree-depth hits, and minimum E-BFMI above 0.66 in both models (Appendix Table~\ref{tab:model-comparison}). Posterior predictive checks reproduced the observed injury counts closely: the full posterior predictive distribution of replicated injured-athlete counts was centered on 115.5 (observed 115) for preparation and 149.5 (observed 149) for competition, with the observed count falling centrally within each replicated distribution (Figure~\ref{fig:posterior-predictive}). Every athlete-level PSIS-LOO Pareto-$k$ value was at most 0.5 in both models, supporting the reliability of the leave-one-out estimates. The exposure model's athlete-grouped ELPD exceeded the baseline model's by 2.22 (SE 2.30) (Appendix Table~\ref{tab:model-comparison}): the exposure terms did not clearly improve held-out predictive performance, even though the competition-period dose contrast itself had a high posterior probability of being positive. These two findings are compatible: a single, spline-parameterized, period-specific exposure effect concentrated in a subset of the competition-period users can produce a credible marginal association without moving whole-model out-of-sample predictive accuracy by a detectable amount in a cohort of this size.

\subsection{Sensitivity analyses}
\label{sec:results-sensitivity}
The direction of the competition-period dose association was stable across both pre-specified sensitivity analyses, but its magnitude was not (Table~\ref{tab:sensitivity}, Figure~\ref{fig:competition-dose-sensitivity}). Under a linear carbon-share parameterization, the association reduced to 6.0 pp (95\% CrI $-0.3$ to $12.6$; posterior probability 96.9\%), with the credible interval marginally including zero. Under wider priors, the association increased to 18.7 pp (95\% CrI 5.2 to 31.9; posterior probability 99.5\%). Taken together, the consistent positive direction across three structurally different specifications supports treating the competition-period dose association as a genuine feature of these data, while the sensitivity of its magnitude to the assumed dose--response shape argues against reporting 14.3 pp as a stable, precise effect size.

\subsection{Population-weighted and subgroup estimates}
\label{sec:results-weighted-subgroups}
After poststratification to national age--sex population strata, the competition-period dose association changed little: 14.9 pp (95\% CrI 3.3 to 26.4) weighted versus 14.3 pp (95\% CrI 2.7 to 25.8) unweighted (Appendix Table~\ref{tab:weighted-vs-unweighted}, Figure~\ref{fig:weighted-dose-by-sex}, Figure~\ref{fig:weighted-dose-by-age}). The two adoption contrasts and the preparation-period dose contrast likewise remained centered near no association after weighting. Conditional on sex or age group, the competition-period dose association retained a consistent point estimate (14.3--15.5 pp) and an essentially unchanged posterior probability of a positive association (99.4\% in every stratum examined; Appendix Table~\ref{tab:subgroups}), providing no descriptive evidence that the association was concentrated in, or specific to, a single sex or age-group stratum.

\begin{figure}[htbp]
\centering
\includegraphics[width=0.85\textwidth]{fig_carbon_share_distribution.png}
\caption{Observed distribution of carbon-plate insole share of reported training-type frequency, among all 352 athletes, by period.}
\label{fig:carbon-share-distribution}
\end{figure}

\begin{figure}[htbp]
\centering
\includegraphics[width=0.95\textwidth]{fig_descriptive_carbon_share_by_stratum.png}
\caption{Observed distribution of carbon-plate insole share of reported training-type frequency, by period, age group, and sex. Curves are boundary-corrected kernel density estimates over the observed discrete-ratio share values (Section~\ref{sec:exposure}), shown to compare distributional overlap across strata; they are a smoothing device for visualization.}
\label{fig:descriptive-carbon-share-by-stratum}
\end{figure}

\begin{figure}[htbp]
\centering
\includegraphics[width=0.85\textwidth]{fig_posterior_predictive.png}
\caption{Posterior predictive check: full posterior distribution of the model-replicated number of injured athletes (4000 posterior draws of the primary model), by period. Dashed vertical line: observed count.}
\label{fig:posterior-predictive}
\end{figure}

\begin{figure}[htbp]
\centering
\includegraphics[width=0.8\textwidth]{fig_posterior_association_distributions.png}
\caption{Full posterior distributions of the adjusted prevalence-difference estimand for the adoption and dose contrasts, by period. Grey band: region of practical equivalence ($\pm$2 pp).}
\label{fig:posterior-associations}
\end{figure}

\begin{figure}[htbp]
\centering
\includegraphics[width=0.85\textwidth]{fig_spline_dose_response.png}
\caption{Fitted carbon-share dose response among users. Curves are conditional injury odds relative to a user with 50\% carbon share; lines show posterior medians, bands 95\% credible intervals, and rugs observed user shares. Gold bands, labeled vertical lines, and outlined points mark the internal spline knots (37.5\% and 57.1\%).}
\label{fig:spline-dose-response}
\end{figure}

\begin{figure}[htbp]
\centering
\includegraphics[width=0.8\textwidth]{fig_competition_dose_sensitivity.png}
\caption{Sensitivity analysis: full posterior distribution of the competition-period dose association (75th vs.\ 25th percentile carbon share among users) under the primary nonlinear specification, a linear share specification, and wider priors.}
\label{fig:competition-dose-sensitivity}
\end{figure}

\begin{figure}[htbp]
\centering
\includegraphics[width=0.75\textwidth]{fig_weighted_dose_by_sex.png}
\caption{Competition-period dose association, population-weighted and stratified by sex.}
\label{fig:weighted-dose-by-sex}
\end{figure}

\begin{figure}[htbp]
\centering
\includegraphics[width=0.75\textwidth]{fig_weighted_dose_by_age.png}
\caption{Competition-period dose association, population-weighted and stratified by age group.}
\label{fig:weighted-dose-by-age}
\end{figure}

\clearpage
\begingroup
\footnotesize
\begin{longtable}{L{3.3cm} L{1.9cm} C{1.7cm} C{1.7cm} C{1.6cm} C{1.6cm}}
\caption{Baseline characteristics of the balanced analysis cohort ($N=352$), by competition sex. Values are mean (SD), median [Q1, Q3], or n (\%). SMD is the standardized mean difference, women vs.\ men (Cohen's $d$ for continuous/ordinal rows, computed on log-odds for proportions); $|\text{SMD}| \geq 0.20$ is conventionally a small imbalance and $\geq 0.50$ a moderate one \citep{table1guidelines}. Population-weighted columns apply the age--sex poststratification weights of Section~\ref{sec:weighting} (Appendix Table~\ref{tab:weight-diagnostics}).}
\label{tab:table1}\\
\toprule
Characteristic & Statistic & Men ($n=194$) & Women ($n=158$) & SMD sample & SMD weighted\\
\midrule
\endfirsthead
\multicolumn{6}{l}{\textit{(continued)}}\\
\toprule
Characteristic & Statistic & Men ($n=194$) & Women ($n=158$) & SMD sample & SMD weighted\\
\midrule
\endhead
\midrule \multicolumn{6}{r}{\textit{(continued on next page)}}\\
\endfoot
\bottomrule
\endlastfoot
\multicolumn{6}{l}{\textbf{Demographics}}\\
Age, years & Mean (SD) & 21.3 (3.1) & 21.9 (3.1) & 0.198 & 0.137 \\
Age group: Under 20 & n (\%) & 71 (36.6) & 43 (27.2) & $-0.202$ & $-0.144$ \\
Age group: Under 23 & n (\%) & 63 (32.5) & 53 (33.5) & 0.023 & 0.060 \\
Age group: Senior & n (\%) & 60 (30.9) & 62 (39.2) & 0.175 & 0.093 \\
\multicolumn{6}{l}{\textbf{Anthropometrics}}\\
Height, cm & Mean (SD) & 180.6 (5.9) & 169.2 (6.8) & $-1.809$ & $-1.749$ \\
Weight, kg & Mean (SD) & 70.5 (7.8) & 57.8 (6.0) & $-1.831$ & $-1.847$ \\
\multicolumn{6}{l}{\textbf{Athletics history}}\\
Experience, years & Mean (SD) & 10.1 (4.4) & 11.0 (3.9) & 0.225 & 0.220 \\
\multicolumn{6}{l}{\textbf{Injury}}\\
Prior injury (previous 2 seasons) & n (\%) & 156 (80.4) & 111 (70.3) & $-0.237$ & $-0.250$ \\
Any current-season injury & n (\%) & 120 (61.9) & 84 (53.2) & $-0.177$ & $-0.210$ \\
Preparation-period injury & n (\%) & 69 (35.6) & 46 (29.1) & $-0.138$ & $-0.165$ \\
Competition-period injury & n (\%) & 88 (45.4) & 61 (38.6) & $-0.137$ & $-0.139$ \\
\multicolumn{6}{l}{\textbf{Carbon-plate insole exposure}}\\
Any preparation carbon use & n (\%) & 159 (82.0) & 128 (81.0) & $-0.024$ & $-0.041$ \\
Any competition carbon use & n (\%) & 173 (89.2) & 135 (85.4) & $-0.112$ & $-0.130$ \\
Preparation carbon share (among users) & Median [Q1,Q3] & 0.4 [0.2, 0.6] & 0.4 [0.3, 0.6] & 0.125 & 0.115 \\
Competition carbon share (among users) & Median [Q1,Q3] & 0.5 [0.3, 0.7] & 0.5 [0.4, 0.8] & 0.184 & 0.161 \\
\multicolumn{6}{l}{\textbf{Training frequency (period total)}}\\
Preparation total type-frequency & Median [Q1,Q3] & 8 [6, 11] & 8 [6, 13] & 0.098 & 0.069 \\
Competition total type-frequency & Median [Q1,Q3] & 7 [5, 10] & 7 [5, 11] & 0.121 & 0.106 \\
\multicolumn{6}{l}{\textbf{Discipline}}\\
Sprint & n (\%) & 64 (33.0) & 43 (27.2) & $-0.126$ & $-0.122$ \\
Middle distance & n (\%) & 38 (19.6) & 31 (19.6) & 0.001 & 0.015 \\
Long distance & n (\%) & 36 (18.6) & 23 (14.6) & $-0.108$ & $-0.064$ \\
Hurdles & n (\%) & 51 (26.3) & 45 (28.5) & 0.049 & 0.023 \\
Horizontal jumps & n (\%) & 18 (9.3) & 26 (16.5) & 0.216 & 0.183 \\
Pole vault & n (\%) & 10 (5.2) & 9 (5.7) & 0.024 & 0.001 \\
Combined events & n (\%) & 16 (8.2) & 25 (15.8) & 0.234 & 0.195 \\
\end{longtable}
\endgroup

\begin{table}[htbp]
\centering
\caption{Sample flow from survey respondents to the balanced analysis cohort.}
\label{tab:sample-flow}
\begin{tabular}{lr}
\toprule
Stage & n \\
\midrule
Consented, eligible, complete survey & 355 \\
\quad Excluded: incomplete injury timing & $-1$ \\
\quad Excluded: zero total eligible frequency in a period & $-2$ \\
Balanced analysis cohort (athletes) & 352 \\
Athlete-period records (preparation + competition) & 704 \\
\midrule
Preparation-period injuries & 115 \\
Competition-period injuries & 149 \\
\bottomrule
\end{tabular}
\end{table}

\begin{table}[htbp]
\centering
\caption{Mean reported carbon share of training-type frequency, overall and by sex and age group. Carbon share is defined in Section~\ref{sec:exposure} and includes non-users (share $=0$).}
\label{tab:descriptive-carbon}
\begin{tabular}{llrrr}
\toprule
Stratum & Level & n & Preparation, mean share & Competition, mean share \\
\midrule
Overall & Overall & 352 & 0.357 & 0.477 \\
Sex & Men & 194 & 0.349 & 0.468 \\
 & Women & 158 & 0.366 & 0.488 \\
Age group & Under 20 & 114 & 0.365 & 0.453 \\
 & Under 23 & 116 & 0.379 & 0.489 \\
 & Senior & 122 & 0.328 & 0.488 \\
\bottomrule
\end{tabular}
\end{table}

\begingroup
\footnotesize
\begin{longtable}{L{3.3cm} C{0.7cm} C{1.15cm} C{0.95cm} C{1.15cm} C{1.15cm} C{1.05cm} C{1.15cm} C{1.05cm} C{1.05cm}}
\caption{Characteristics and injury prevalence by specialty group. Groups collapse the seven athletics disciplines (Appendix Table~\ref{tab:discipline-support}) into four physiologically motivated categories: Sprint/hurdles (sprint or hurdles), Endurance (middle or long distance), Jumps (horizontal jumps or pole vault), and Combined events. Because disciplines and, consequently, groups are multi-select, 26 of 352 athletes belong to more than one group.}
\label{tab:specialty-group}\\
\toprule
& $n$ & Age, y & Women, \% & Height, cm & Weight, kg & Exp., y & Prior inj., \% & Any inj., \% & Comp.\ inj., \% \\
\midrule
\endfirsthead
\multicolumn{10}{l}{\textit{(continued)}}\\
\toprule
& $n$ & Age, y & Women, \% & Height, cm & Weight, kg & Exp., y & Prior inj., \% & Any inj., \% & Comp.\ inj., \% \\
\midrule
\endhead
\midrule \multicolumn{10}{r}{\textit{(continued on next page)}}\\
\endfoot
\bottomrule
\endlastfoot
Sprint/hurdles & 169 & 21.4 (3.0) & 42.6 & 175.8 (8.7) & 66.9 (9.3) & 10.0 (3.9) & 75.7 & 60.9 & 44.4 \\
Endurance & 109 & 21.6 (3.2) & 42.2 & 174.0 (8.2) & 59.6 (7.9) & 10.3 (4.3) & 67.9 & 44.0 & 30.3 \\
Jumps & 62 & 22.2 (3.5) & 54.8 & 175.6 (9.0) & 65.8 (9.6) & 11.5 (4.6) & 83.9 & 64.5 & 51.6 \\
Combined events & 41 & 22.0 (3.0) & 61.0 & 176.6 (7.4) & 67.8 (8.3) & 11.9 (4.4) & 82.9 & 70.7 & 48.8 \\
\end{longtable}
\endgroup

\begin{table}[htbp]
\centering
\footnotesize
\caption{Carbon-plate insole use and training-type frequency by specialty group (grouping as in Table~\ref{tab:specialty-group}).}
\label{tab:specialty-group-carbon}
\begin{tabular}{lcccc}
\toprule
Specialty group & Prep.\ carbon share & Comp.\ carbon share & Prep.\ total freq., Mdn & Comp.\ total freq., Mdn \\
\midrule
Sprint/hurdles & 0.39 & 0.56 & 8 & 7 \\
Endurance & 0.31 & 0.34 & 10 & 9 \\
Jumps & 0.33 & 0.47 & 6 & 5 \\
Combined events & 0.41 & 0.52 & 9 & 9 \\
\bottomrule
\end{tabular}
\end{table}

\begin{table}[htbp]
\centering
\footnotesize
\caption{Anatomical location of reported injuries (259 injuries among 204 injured athletes; multi-select item, so location counts sum to more than 259).}
\label{tab:injury-location}
\begin{tabular}{lrrr}
\toprule
Location & Injuries, $n$ & \% of injuries & \% of injured athletes ($n=204$) \\
\midrule
Thigh (anterior or posterior) & 97 & 37.5 & 39.2 \\
Lower leg (anterior or posterior) & 65 & 25.1 & 27.9 \\
Foot (anterior or posterior) & 50 & 19.3 & 22.1 \\
Ankle (medial or lateral) & 36 & 13.9 & 15.2 \\
Groin & 26 & 10.0 & 12.3 \\
Hip & 14 & 5.4 & 6.9 \\
Head, upper limb, or trunk & 11 & 4.2 & 5.4 \\
\bottomrule
\end{tabular}
\end{table}

\begin{table}[htbp]
\centering
\footnotesize
\caption{Type of reported lower-limb injuries. The type item was asked only for injuries localized to the hip, groin, thigh, lower leg, ankle, or foot (248 of 259 injuries; the questionnaire ended after the location item for the 11 head/upper-limb/trunk injuries); multi-select item, so counts sum to more than 248.}
\label{tab:injury-type}
\begin{tabular}{lrrr}
\toprule
Type & Injuries, $n$ & \% of lower-limb injuries & \% of injured athletes ($n=204$) \\
\midrule
Tendinopathy & 71 & 28.6 & 29.9 \\
Bone stress injury or stress fracture & 65 & 26.2 & 25.0 \\
Muscle strain & 55 & 22.2 & 26.5 \\
Other & 39 & 15.7 & 16.7 \\
Ligament injury & 18 & 7.3 & 8.8 \\
Plantar fasciopathy & 16 & 6.5 & 7.8 \\
Joint sprain & 14 & 5.6 & 6.9 \\
Tendon rupture & 10 & 4.0 & 3.9 \\
\bottomrule
\end{tabular}
\end{table}

\begin{table}[htbp]
\centering
\footnotesize
\caption{Anatomical location of the three leading injury types (\% of injuries of that type; multi-select, so rows sum to more than 100\%). Location categories with 0\% in all three rows (hip, groin, head/upper limb/trunk) are omitted.}
\label{tab:injury-type-by-location}
\begin{tabular}{lrrrr}
\toprule
Type ($n$ injuries) & Thigh & Lower leg & Ankle & Foot \\
\midrule
Tendinopathy (71) & 22.5 & 33.8 & 22.5 & 35.2 \\
Bone stress injury/stress fracture (65) & 58.5 & 26.2 & 7.7 & 15.4 \\
Muscle strain (55) & 74.5 & 27.3 & 3.6 & 0.0 \\
\bottomrule
\end{tabular}
\end{table}

\begin{table}[htbp]
\centering
\caption{Primary adjusted associations between carbon-plate insole exposure and injury prevalence, by period and contrast. Adjusted prevalence difference and ratio are posterior-standardized (Section~\ref{sec:estimands}); credible intervals are 95\% equal-tailed. Adoption compares the median-share user with a non-user; dose compares the 75th with the 25th percentile of share among users.}
\label{tab:primary-associations}
\begin{tabular}{llrrr}
\toprule
Period & Contrast & Difference, pp (95\% CrI) & Ratio (95\% CrI) & P(positive) \\
\midrule
Preparation & Adoption & 5.4 ($-6.3$ to $16.2$) & 1.21 (0.84 to 1.74) & 82.0\% \\
Preparation & Dose & $-0.8$ ($-10.6$ to $9.1$) & 0.99 (0.73 to 1.30) & 43.8\% \\
Competition & Adoption & 4.8 ($-8.7$ to $17.7$) & 1.16 (0.82 to 1.65) & 77.1\% \\
Competition & Dose & \textbf{14.3 (2.7 to 25.8)} & \textbf{1.41 (1.06 to 1.86)} & \textbf{99.0\%} \\
\bottomrule
\end{tabular}
\end{table}

\begin{table}[htbp]
\centering
\caption{Sensitivity of the competition-period dose association (75th vs.\ 25th percentile share among users) to model specification.}
\label{tab:sensitivity}
\begin{tabular}{lrr}
\toprule
Model & Difference, pp (95\% CrI) & P(positive) \\
\midrule
Primary (nonlinear spline share) & 14.3 (2.7 to 25.8) & 99.0\% \\
Linear share & 6.0 ($-0.3$ to $12.6$) & 96.9\% \\
Wider priors & 18.7 (5.2 to 31.9) & 99.5\% \\
\bottomrule
\end{tabular}
\end{table}

\clearpage
